\begin{abstract}
Every nonexpansive selfmap of a nonempty closed bounded convex subset of a real reflexive Banach space has a fixed point. This resolves the reflexive-space fixed point problem for the given norm.
\end{abstract}
\section{Introduction}
\label{sec:introduction}

A Banach space has the \emph{fixed point property for nonexpansive maps}
if every nonexpansive selfmap of each nonempty closed bounded convex
subset has a fixed point. Here nonexpansiveness is measured in the given
norm. We prove the following theorem.

\begin{theorem}\label{thm:main}
Let $X$ be a real reflexive Banach space, and let $C\subseteq X$ be
nonempty, norm closed, bounded, and convex. If $F:C\to C$ satisfies
\[
  \norm{F(a)-F(b)}\le \norm{a-b}\qquad(a,b\in C),
\]
then $F$ has a fixed point.
\end{theorem}

Theorem~\ref{thm:main} establishes the positive assertion in the
reflexive-space fixed point problem for the given norm.

Bruck's theorem for commuting families gives a further consequence.

\begin{corollary}[Commuting families]\label{cor:commuting}
Let $X$ be a real reflexive Banach space, and let $C\subseteq X$ be
nonempty, norm closed, bounded, and convex. For any pairwise commuting
family $\mathcal S$ of nonexpansive selfmaps of $C$, the common fixed set
\[
  \operatorname{Fix}(\mathcal S)
  =\{x\in C:T(x)=x\text{ for every }T\in\mathcal S\}
\]
is nonempty and is a nonexpansive retract of $C$. That is, there is a
nonexpansive map $R:C\to\operatorname{Fix}(\mathcal S)$ that is the
identity on $\operatorname{Fix}(\mathcal S)$.
\end{corollary}

\begin{proof}
For the empty family, take $R=\operatorname{id}_C$. Otherwise,
Theorem~\ref{thm:main} applies to every nonempty closed bounded convex
subset $D$ of $C$, since a relatively norm-closed subset of the
norm-closed set $C$ is norm closed in $X$. Thus $C$
